% ============================================================================
% CONSISTENCY GUIDE: Changes needed in Chapters 1-3
% to match the conventions used in Chapters 4-6
% ============================================================================
%
% This document lists every change needed in the initial draft (Chapters 1-3)
% so that the full thesis reads as one continuous, consistent document.
%
% Priority levels:
%   [MUST]   = Will cause confusion or inconsistency if not fixed
%   [SHOULD] = Improves quality and consistency
%   [NICE]   = Minor polish, optional
%
% ============================================================================

% ============================================================================
% 1. CONTRACTIONS — Remove ALL contractions
% Priority: [MUST]
% ============================================================================
%
% Ch4-6 use ZERO contractions. Ch1-3 uses many.
% A professor will notice the sudden shift from informal to formal.
%
% Search and replace in Ch1-3:
%   "isn't"     → "is not"
%   "don't"     → "do not"
%   "can't"     → "cannot"
%   "won't"     → "will not"
%   "didn't"    → "did not"
%   "doesn't"   → "does not"
%   "shouldn't" → "should not"
%   "wouldn't"  → "would not"
%   "it's"      → "it is" (when meaning "it is") or "its" (when possessive)
%   "there's"   → "there is"
%   "that's"    → "that is"
%   "we've"     → "we have"
%   "they're"   → "they are"
%   "we're"     → "we are"
%   "haven't"   → "have not"
%   "wasn't"    → "was not"
%   "weren't"   → "were not"
%   "couldn't"  → "could not"
%
% Exception: Keep contractions inside code listings (lstlisting blocks)

% ============================================================================
% 2. DASHES — Remove ALL em-dashes and en-dashes from text
% Priority: [MUST]
% ============================================================================
%
% Ch4-6 use ZERO dashes in running text.
% Ch1-3 may have em-dashes (---) or en-dashes (--) in prose.
%
% Replace:
%   "word --- word" → "word, word" (use comma instead of em-dash)
%   "word -- word"  → "word, word" (use comma instead of en-dash)
%   "X--Y" (number ranges) → "X to Y" (write out "to")
%
% Keep dashes ONLY in:
%   - Code listings (e.g., --flag, --option)
%   - URLs
%   - LaTeX commands

% ============================================================================
% 3. SEMICOLONS AND COLONS — Replace with commas
% Priority: [MUST]
% ============================================================================
%
% Ch4-6 use ZERO semicolons or mid-sentence colons in running text.
% Ch1-3 likely uses some.
%
% Replace:
%   "sentence; sentence" → "sentence, sentence" (or restructure)
%   "concept: explanation" → "concept, explanation" (mid-sentence colons)
%
% Keep colons ONLY in:
%   - Section/subsection titles (e.g., "Phase 1: Local Development")
%   - Bold labels (e.g., \textbf{RQ1}: ...)
%   - List item labels (e.g., \item \textbf{Tool}: Description)
%   - Code listings
%   - URLs

% ============================================================================
% 4. AI-FLAGGED WORDS — Remove or replace
% Priority: [MUST]
% ============================================================================
%
% Ch4-6 have ZERO instances of these words.
% Search Ch1-3 and replace any found:
%
%   "straightforward" → "simple" or "direct" or remove
%   "ensure"          → "verify" or "confirm" or "check"
%   "significant"     → use sparingly (max 1-2 per chapter)
%   "substantial"     → "large" or "extensive" or "wide"
%   "comprehensive"   → "thorough" or "complete" or "full"
%   "fundamental"     → "core" or "basic" or "central"
%   "leverage"        → "use" or "apply"
%   "utilize"         → "use"
%   "robust"          → "reliable" or "strong"
%   "paradigm"        → "approach" or "model"
%   "delve"           → "examine" or "explore"
%   "interestingly"   → remove or rephrase
%   "notably"         → remove or rephrase
%   "importantly"     → remove or rephrase
%   "furthermore"     → remove or start new sentence
%   "moreover"        → remove or start new sentence
%   "additionally"    → remove or start new sentence
%   "optimal"         → "best" or "right" or "suitable"
%   "intricate"       → "complex" or "detailed"
%
% Also avoid:
%   "it is worth noting that" → remove, state the fact directly
%   "it is important to note" → remove, state the fact directly
%   "in order to"             → "to"

% ============================================================================
% 5. MODEL NAMES — Use consistent naming
% Priority: [MUST]
% ============================================================================
%
% Ch4-6 use these exact model names. Update Ch1-3 to match:
%
%   Gemma 3 27B           (not "Gemma 27B" or "Gemma-3-27B")
%   Qwen 2.5-Coder 32B   (not "Qwen 32B" or "Qwen2.5-Coder-32B")
%   Phi-4 14B             (not "Phi 14B" or "Phi-4-14B" or "Phi 4 3.8B")
%   DeepSeek-Coder-V2-Lite 16B
%   DeepSeek-Coder 33B
%   DeepSeek R1
%   Yi 34B
%   Mixtral 8x7B 46.7B   (mention "46.7B total, 12.9B active parameters")
%   WizardCoder 34B
%   CodeLlama 34B
%   Qwen 2.5-Coder 14B
%   Qwen 2.5-Coder 7B
%   Qwen 2.5-Coder 1.5B
%   Gemma 3 4B
%   Gemma 3 9B
%
% If Ch1-3 mentions model counts, it should say "14 models"

% ============================================================================
% 6. TOOL NAMES — Use consistent naming with citations
% Priority: [SHOULD]
% ============================================================================
%
% Ch4-6 cite all tools on first mention. Ch1-3 should do the same:
%
%   AFL++ \cite{AFLPlusPlus}
%   libFuzzer \cite{LibFuzzer2015}
%   cifuzz \cite{cifuzz}
%   Ollama \cite{ollama}
%   Azure OpenAI \cite{azure_openai}
%   Buildah \cite{buildah}
%   AddressSanitizer (ASan)
%   UndefinedBehaviorSanitizer (UBSan)
%
% Expand abbreviations on first use:
%   "ASan" → "AddressSanitizer (ASan)" on first mention
%   "UBSan" → "UndefinedBehaviorSanitizer (UBSan)" on first mention
%   "AUTOSAR" → "AUTomotive Open System ARchitecture (AUTOSAR)" on first mention
%   "MISRA" → "Motor Industry Software Reliability Association (MISRA)" on first mention

% ============================================================================
% 7. NUMBERS AND PERCENTAGES — Use consistent formatting
% Priority: [SHOULD]
% ============================================================================
%
% Ch4-6 conventions:
%   - Coverage percentages: one decimal place (45.06\%, 43.08\%, 34.26\%)
%   - Cost figures: Euro symbol with comma thousands (€1,452)
%   - Token counts: comma-separated (43,100 tokens) or with "k" (43.1k)
%   - Time: written out (37 minutes, 1 hour 12 minutes)
%   - Model count: "14" (numeral, not "fourteen")
%   - Small numbers in text: use numerals for technical values, words for general
%   - Number ranges: "X to Y" (not "X-Y" or "X--Y")
%
% Update Ch1-3 to follow same conventions

% ============================================================================
% 8. LaTeX FORMATTING — Use consistent commands
% Priority: [SHOULD]
% ============================================================================
%
% Ch4-6 conventions:
%   - Non-breaking spaces: Table~\ref{}, Figure~\ref{}, Section~\ref{}
%     (never "Table \ref{}" with a regular space)
%   - Memory sizes: 20~GB, 16~GB (tilde for non-breaking space)
%   - Code inline: \texttt{code} for file names, commands, paths
%   - Emphasis: \textbf{} for key terms on first introduction
%   - Lists: \begin{itemize} with \item entries
%   - Tables: \begin{table}[htbp] with \centering
%   - All tables use either \small or \resizebox{\textwidth}{!}{}
%     (use \small for tables with fewer than 6 columns,
%      use \resizebox for tables with 6+ columns or long text)

% ============================================================================
% 9. CITATION KEYS — Use matching keys from literature.bib
% Priority: [MUST]
% ============================================================================
%
% If Ch1-3 has any \cite{} commands, verify the keys match literature.bib.
% Current valid keys in literature.bib:
%
%   Miller1990, AFL, AFLPlusPlus, LibFuzzer2015, SAGE2008,
%   DeepXplore2017, NeuFuzz2019, MLFuzzReview2020, FuzzSurvey2018,
%   ChatAFL2024, HGFuzzer2024, TitanFuzz2023, FuzzGPT2024,
%   LLM4Fuzz2024, CodaMosa2023, SBFT2024, PromptFuzz2024,
%   ISO21434, ISO26262, UNECE_R155, AUTOSAR2022, FuzzBench2021,
%   ContinuousDelivery2010, SandiaFuzzing2007, DevSecOps2020,
%   Broy2006, DESTINA2024, cifuzz, autosar, azure_openai,
%   azure_private_link, ollama, gpt4, llm4fuzz, efficient_finetuning,
%   buildah
%
% Any citation key in Ch1-3 not in this list needs to be either:
%   (a) added to literature.bib, or
%   (b) changed to match an existing key

% ============================================================================
% 10. CONTROVERSIAL LANGUAGE — Check and neutralize
% Priority: [MUST]
% ============================================================================
%
% Ch4-6 follow strict rules about controversial content.
% Apply same rules to Ch1-3:
%
% DO NOT use:
%   - Language criticizing CARIAD (no "barriers", "obstacles", "limitations")
%   - Language criticizing any tool (no "failed", "broken", "buggy")
%   - Salary or compensation data
%   - Comparisons that imply one company/tool is inferior
%   - "conservative" when describing processes or cultures
%   - "bureaucratic" or "slow" about approval processes
%
% DO use:
%   - Neutral framing: "required additional configuration"
%   - Factual descriptions: "the process took one month"
%   - Positive framing: "security policies that govern..."
%   - "did not produce" instead of "failed to produce"
%   - "encountered" instead of "struggled with"

% ============================================================================
% 11. PHASE NAMES — Align with Ch4 implementation phases
% Priority: [SHOULD]
% ============================================================================
%
% Ch3 methodology defines phases as:
%   "Phase 1: Baseline Assessment"
%   "Phase 2: Model Evaluation"
%   "Phase 3: Fine-tuning"
%
% Ch4 implements these as:
%   "Phase 1: Local Development Environment"
%   "Phase 2: Cloud Integration with Azure OpenAI"
%   "Phase 3: LoRA Fine-tuning"
%
% Ch4 already has a bridging sentence explaining the mapping.
% If Ch3 can be edited, consider adding a forward reference:
%   "These methodology phases are implemented in Chapter 4 as
%    local development, cloud integration, and LoRA fine-tuning."
%
% Alternatively, just make sure Ch3's phase descriptions are
% compatible with what Ch4 actually describes.

% ============================================================================
% 12. LIBRARY NAMES — Use consistent naming
% Priority: [SHOULD]
% ============================================================================
%
% Ch4-6 use these exact library names:
%   yaml-cpp (primary target)
%   fmt
%   pugixml
%   jsoncons
%   glm
%   spdlog
%
% Ch3 mentions "25 open-source C/C++ libraries"
% If Ch3 lists specific libraries, make sure names match exactly.

% ============================================================================
% 13. KEY NUMBERS TO KEEP CONSISTENT
% Priority: [MUST]
% ============================================================================
%
% These numbers appear in Ch4-6 and must match Ch1-3:
%
%   14 models evaluated
%   6 target libraries (yaml-cpp, fmt, pugixml, jsoncons, glm, spdlog)
%   25 open-source C/C++ repositories (initial corpus)
%   709 prompt-response pairs (Qwen fine-tuning dataset)
%   172 prompt-response pairs (Gemma fine-tuning dataset)
%   45.06\% coverage (Gemma 3 27B on yaml-cpp)
%   43.08\% coverage (Qwen 2.5-Coder 32B on yaml-cpp)
%   34.26\% coverage (Phi-4 14B on yaml-cpp)
%   33\% coverage improvement from fine-tuning (Qwen 32B)
%   55\% coverage improvement from fine-tuning (Gemma 27B)
%   5 months project duration (May to September)
%   1 month for Azure Private Link setup
%   Apple M1 Pro with 16 GB unified memory
%   €0.20 per 1,000 tokens (Azure OpenAI pricing)
%   Cost tiers: Light €74, Moderate €219, Heavy €726, Enterprise €1,452
%
% If Ch1 abstract mentions any of these numbers, they must match exactly.

% ============================================================================
% 14. SENTENCE PATTERNS TO AVOID
% Priority: [NICE]
% ============================================================================
%
% Ch4-6 avoid these patterns. Apply to Ch1-3 as well:
%
%   - Sentences starting with "It is" (passive, AI-like)
%   - Multiple consecutive sentences starting with "The"
%   - Sentences starting with "This is" without a clear referent
%   - Overly long sentences (break at 40+ words)
%   - One-sentence paragraphs (combine with neighbors or expand)
%
% Preferred patterns (matching Ch4-6 style):
%   - Start with subject: "We evaluated...", "The model produced..."
%   - Vary sentence openers: "After...", "During...", "To achieve..."
%   - Use first person plural: "We found...", "Our results show..."

% ============================================================================
% SUMMARY CHECKLIST
% ============================================================================
%
% Before submitting Ch1-3, verify:
%
% [ ] Zero contractions in running text
% [ ] Zero em-dashes (---) or en-dashes (--) in running text
% [ ] Zero semicolons in running text
% [ ] Zero mid-sentence colons (keep only in labels/titles)
% [ ] Zero AI-flagged words (straightforward, ensure, leverage, etc.)
% [ ] All model names match Ch4-6 naming exactly
% [ ] All tool names cited on first mention
% [ ] All citation keys present in literature.bib
% [ ] All numbers match Ch4-6 exactly (14 models, 45.06\%, etc.)
% [ ] No controversial language about CARIAD or any tools
% [ ] Phase names compatible with Ch4 implementation
% [ ] Figure/Table references use non-breaking space (~)
% [ ] Consistent "we" usage throughout

\endinput
